\documentclass[10pt,a4paper]{amsart}
\usepackage[margin=19mm]{geometry}
\usepackage{amsmath,amssymb,mathtools}
\usepackage[scr=rsfs,cal=euler]{mathalfa}
\usepackage{xfrac}
\usepackage[unicode,hidelinks,bookmarks=false]{hyperref}
\newcommand{\ie}{i.e.\ }
\newcommand{\cf}{cf.\ }
\newcommand{\resp}{respectively\ }
\newcommand{\Subsection}{\subsection}
\providecommand{\nobreakdash}{\nobreak\hbox{-}}
\setlength{\emergencystretch}{2em}
\multlinegap0pt
\usepackage{bm}
\usepackage{bbm}
\usepackage{dsfont}
\usepackage{xspace}
\usepackage{cancel}
\usepackage{mathtools}
\usepackage{amsthm}
\makeatletter
\@ifundefined{theorem}{\newtheorem{theorem}{Theorem}}{}
\@ifundefined{lemma}{\newtheorem{lemma}[theorem]{Lemma}}{}
\makeatother
\newcommand{\R}{\mathbb{R}}
\newcommand{\one}{\mathbf{1}}
\newcommand{\vhn}{\hat{V}_{\mathrm{ney}}}
\newcommand{\xbs}{\bar{x}}
\newcommand{\ybh}{\hat{\bar{Y}}}
\newcommand{\yb}{\bar{Y}}
\title[Randomization Inference with Concentration Inequalities]{Randomization Inference with Concentration Inequalities}
\author{Tobias Freidling}
\date{Source: arXiv:2609.18586v1 (CC BY 4.0)}
\begin{document}
\maketitle

\noindent\emph{Source and licence.} Randomization Inference with Concentration Inequalities. Version arXiv:2609.18586v1 (CC BY 4.0), distributed under the Creative Commons Attribution 4.0 International licence, as recorded in the arXiv licence metadata for this exact version and shown on the abstract page https://arxiv.org/abs/2609.18586v1. Full rights notice: licenses/arxiv-2609.18586-CC-BY-4.0.txt.\par\smallskip

\noindent\textit{Editorial scope.} Counted unit: the Lemma whose source label is \texttt{lem:self-bounded} in the article (source0.tex lines 1040--1053, source label \texttt{lem:self-bounded}), delivered together with the article's own complete original proof of it (source0.tex lines 1054--1131, one \texttt{proof} environment of 78 lines). No mathematics is written, repaired or re-derived: statement and proof are the article's own text line for line. Required context retained uncounted: lem:emp-mean-perturb (lines 1505--1512). Every definition, display and label of those lines is retained unchanged. Omitted from this package and not redistributed: 1--1039, 1132--1504, 1513--1536. Nothing inside a retained excerpt is omitted or altered: every retained body is delivered exactly as the article's lines are written, so no omission receipt of any kind accompanies this package. Citations: the retained excerpts carry no citation command at all, so no bibliography entry of the article is transcribed and no reference is redistributed. Reference adaptation: none was needed and none was made; every reference inside the delivered unit resolves inside it, so no unresolved reference or citation is produced. Printed slips kept literal: no printed word, symbol, hypothesis or number of the counted statement or of its proof was altered. Layout and class: the article's own class is \texttt{imsart2} with packages including amsthm, amsmath, amsfonts, amssymb, natbib, hyperref, graphicx; the wrapper is the lane's pinned \texttt{amsart} editorial wrapper at 10pt on A4, which restates the theorem-like environments the retained text needs and adds the article's own macro definitions that the retained text needs (\texttt{\textbackslash R}, \texttt{\textbackslash one}, \texttt{\textbackslash vhn}, \texttt{\textbackslash xbs}, \texttt{\textbackslash yb}, \texttt{\textbackslash ybh}), with extra wrapper packages (bm, bbm, dsfont, xspace, cancel, mathtools). The delivered numbering is the unit's own. Assets: no retained span contains a figure environment, a graphics-inclusion command, a PDF-page inclusion, a table or a TikZ picture; no image file is copied into this package, the package declares no asset dependency and no image file is redistributed with it. Licence: version arXiv:2609.18586v1 of this article is distributed under the Creative Commons Attribution 4.0 International licence, as recorded in the arXiv licence metadata for this exact version and shown on the abstract page https://arxiv.org/abs/2609.18586v1. A full rights notice is delivered at licenses/arxiv-2609.18586-CC-BY-4.0.txt. Characters: every retained excerpt is pure ASCII, so the delivered engine, XeLaTeX with its default Latin Modern text font, reports no missing character.
\begin{lemma}\label{lem:emp-mean-perturb}
    Let $I$ be a finite index set with $\lvert I \rvert=n$ and let the collection $(x_i)_{i\in I}$ be real-valued. Then, for any $l \in I$,
	%For any $x_1,\ldots,x_n \in \R$ and $j \in \{1,\ldots,n\}$,
	\begin{equation*}
		\sum_{i \in I} (x_i - \xbs)^2 = \sum_{i \in I \setminus\{l\}} (x_i - \xbs_{-l})^2 + \frac{n-1}{n} (x_l - \xbs_{-l})^2,
	\end{equation*}
	where $\xbs:=\frac{1}{n}\sum_{i\in I} x_i$ and $\xbs_{-l}:=\frac{1}{n-1}\sum_{i \in I \setminus\{l\}}^n x_i$.
\end{lemma}

\begin{lemma}[Self-boundedness of $\vhn$]\label{lem:self-bounded}
	Define the swap operator $\tau_{ij}\colon \{0,1\}^n \to \{0,1\}^n$ as
	%\begin{equation*}
    		$\tau_{ij}(z) = (z_1,\ldots, z_{i-1},z_j,z_{i+1},\ldots, z_{j-1}, z_i, z_{j+1},\ldots,z_n)$,
    	%\end{equation*}
	where $i,j\in\{1,\ldots,n\}$. Neyman's variance estimator is self-bounded in the sense that
	\begin{equation*}
		\sum_{i,j=1}^n \big(\vhn(Z)-\vhn(\tau_{ij}Z)\big)_+^2\, \leq\, 2 (b-a)^2\,c(n_0,n_1)\, \vhn(Z),
	\end{equation*}
	where the constant is given by
	\begin{equation*}
		c(n_0,n_1) := \frac{n\,(n_0^2+n_1^2)^2}{n_0n_1^3(n_0-1) + n_1 n_0^3(n_1-1)}.% = \mathcal{O}(\tfrac{1}{n}).
	\end{equation*}
\end{lemma}

\begin{proof}
	We abbreviate the left-hand side of the inequality as $L(Z) := \sum_{i,j=1}^n \big(\vhn(Z)-\vhn(\tau_{ij}Z)\big)_+^2$. First, we notice that only index pairs $(i,j)$ with $Z \neq \tau_{ij} Z$ contribute to the sum; hence,
	\begin{align*}
		L(Z) &= \sum_{i,j=1}^n \big(\vhn(Z)-\vhn(\tau_{ij}Z)\big)_+^2 \one \{Z \neq \tau_{ij} Z\}\\
		&= \sum_{i,j=1}^n \big(\vhn(Z)-\vhn(\tau_{ij}Z)\big)_+^2 \Big(Z_i (1-Z_j) + (1-Z_i)Z_j\Big)\\
		&= 2\sum_{i,j=1}^n Z_i (1-Z_j) \big(\vhn(Z)-\vhn(\tau_{ij}Z)\big)_+^2,
	\end{align*}
	where the last step follows from the symmetry in the summation. Next, we separate the Neyman variance estimator into its two summands which are given by
    \begin{equation*}
        \vhn = \vhn^1+\vhn^0 := \frac{n}{n_1(n_1-1)}\sum_{i \in I_1} (Y_i(1)-\ybh(1))^2 + \frac{n}{n_0(n_0-1)}\sum_{i \in I_0} (Y_i(0)-\ybh(0))^2,
    \end{equation*}
    where $I_k := \{i\colon Z_i=k\}$ for $k\in \{0,1\}$. We fix indices $i\in I_1$ and $j \in I_0$ and consider the first summand under $Z$ and the swapped treatment $\tau_{ij}Z$:
    \begin{align*}
        \vhn^1(Z) &= \frac{n}{n_1(n_1-1)}\sum_{k \in I_1} (Y_k(1)-\ybh(1))^2,\\
        \vhn^1(\tau_{ij}Z) &= \frac{n}{n_1(n_1-1)}\sum_{k \in I_1 \setminus \{i\} \cup \{j\}} (Y_k(1)-\ybh_{-i+j}(1))^2.
    \end{align*}
    Here, $\ybh(1)$ denotes the average of the $Y(1)$-potential outcomes over the indices $I_1$ and $\ybh_{-i+j}(1)$ is the average over the indices $I_1\setminus\{i\} \cup \{j\}$. Applying Lemma~\ref{lem:emp-mean-perturb} with $I = I_1$ and $l=i$ to the first expression and with $I=I_1 \setminus \{i\} \cup \{j\}$ and $l=j$ to the second, we obtain
    \begin{align*}
        \vhn^1(Z) &= \frac{n}{n_1(n_1-1)}\left[\sum_{l' \in I_1 \setminus\{i\}} (Y_{l'}(1)-\ybh_{-i}(1))^2 + \frac{n_1-1}{n_1}(Y_i(1)-\ybh_{-i}(1))^2\right],\\
        \vhn^1(\tau_{ij}Z) &= \frac{n}{n_1(n_1-1)} \left[\sum_{l' \in I_1 \setminus \{i\} \cup \{j\}} (Y_{l'}(1)-\ybh_{-i}(1))^2 + \frac{n_1-1}{n_1} (Y_j(1)-\ybh_{-i}(1))^2\right].
    \end{align*}
    We take the difference and get
    \begin{equation*}
        \vhn^1(Z) - \vhn^1(\tau_{ij}Z) = \frac{n}{n_1^2}\Big[(Y_i(1)-\ybh_{-i}(1))^2-(Y_j(1)-\ybh_{-i}(1))^2\Big].
    \end{equation*}
    %\begin{equation*}
    %    \vhn^1(Z) - \vhn^1(\tau_{ij}Z) = \frac{(Y_i(1)-\ybh_{-i}(1))^2-(Y_j(1)-\ybh_{-i}(1))^2}{n_1^2}.
    %\end{equation*}
    Following the same steps for $\vhn^0$, we obtain
    \begin{equation*}
        \vhn^0(Z) - \vhn^0(\tau_{ij}Z) = \frac{n}{n_0^2}\Big[(Y_j(0)-\ybh_{-j}(0))^2-(Y_i(0)-\ybh_{-j}(0))^2\Big].
    \end{equation*}
    %\begin{equation*}
    %    \vhn^0(Z) - \vhn^0(\tau_{ij}Z) = \frac{(Y_j(0)-\ybh_{-j}(0))^2-(Y_i(0)-\ybh_{-j}(0))^2}{n_0^2}.
    %\end{equation*}

    We can now return to $L(Z)$, plug in our findings above and estimate the resulting expression using the monotonicity of the positive part:
    \begin{align*}
        L(Z) &= \begin{aligned}[t]
            2 n^2 \sum_{i,j=1}^n Z_i (1-Z_j) \bigg(&\frac{(Y_i(1)-\ybh_{-i}(1))^2-(Y_j(1)-\ybh_{-i}(1))^2}{n_1^2}\\
            &+ \frac{(Y_j(0)-\ybh_{-j}(0))^2-(Y_i(0)-\ybh_{-j}(0))^2}{n_0^2}\bigg)_+^2
        \end{aligned}\\[1ex]
        &\leq 2 n^2 \sum_{i,j=1}^n Z_i (1-Z_j) \left(\frac{(Y_i(1)-\ybh_{-i}(1))^2}{n_1^2}
            + \frac{(Y_j(0)-\ybh_{-j}(0))^2}{n_0^2}\right)^2.\\
        %&= 2 \sum_{i,j=1}^n Z_i (1-Z_j) \left(\frac{(Y_i(1)-\yb(1))^2}{(n_1-1)^2}
        %    + \frac{(Y_j(0)-\yb(0))^2}{(n_0-1)^2}\right)^2
    \end{align*}
     %The inequality holds as the positive part is monotonically increasing.
     
     Next, we expand the square and find an upper bound using that $\lvert Y_i(k) - \yb_{-i}(k) \rvert \leq b-a$ for all $i \in \{1,\ldots,n\}$ and $k\in \{0,1\}$. Moreover, we can replace the ``leave-one-out'' averages with their usual counter-parts via $\ybh_{-i}(1)= \frac{n_1\ybh(1)-Y_i(1)}{n_1-1}$ and $\ybh_{-j}(0)= \frac{n_0\ybh(0)-Y_j(0)}{n_0-1}$ and subsequently gather terms:
     {\allowdisplaybreaks
    \begin{align*}
        L(Z) &= \begin{aligned}[t]
        \frac{2 n_0 n^2}{n_1^4} &\sum_{i \in I_1} (Y_i(1)-\ybh_{-i}(1))^4 + \frac{2 n_1 n^2}{n_0^4} \sum_{j \in I_0} (Y_j(0)-\ybh_{-j}(0))^4 \\
        &+ \frac{4n^2}{n_1^2 n_0^2} \sum_{i \in I_1} \sum_{j \in I_0} (Y_i(1)-\ybh_{-i}(1))^2 (Y_j(0)-\ybh_{-j}(0))^2
        \end{aligned}\\
        &\leq\begin{aligned}[t]
        2 (b-a)^2 n^2\Bigg[
        \frac{n_0}{n_1^4} &\sum_{i \in I_1} (Y_i(1)-\ybh_{-i}(1))^2 + \frac{n_1}{n_0^4} \sum_{j \in I_0} (Y_j(0)-\ybh_{-j}(0))^2 \\
        &+ \frac{2}{n_1^2 n_0^2} \sum_{i \in I_1} \sum_{j \in I_0} \lambda (Y_i(1)-\ybh_{-i}(1))^2 + (1-\lambda) (Y_j(0)-\ybh_{-j}(0))^2\Bigg]
        \end{aligned}\\
        &=\begin{aligned}[t]
        2 (b-a)^2 n^2\Bigg[
        &\frac{n_0}{n_1^2(n_1-1)^2} \sum_{i \in I_1} (Y_i(1)-\ybh(1))^2 + \frac{n_1}{n_0^2(n_0-1)^2} \sum_{j \in I_0} (Y_j(0)-\ybh(0))^2 \\
        & + \frac{2 \lambda}{(n_1-1)^2n_0} \sum_{i \in I_1} Y_i(1)-\ybh(1))^2 + \frac{2 (1-\lambda)}{(n_0-1)^2n_1} \sum_{j \in I_0} Y_j(0)-\ybh(0))^2\Bigg]
        \end{aligned}\\
        &=2(b-a)^2 n \left[\left(\frac{n_0}{n_1(n_1-1)} + \frac{2\lambda n_1}{n_0 (n_1-1)}\right) \vhn^1 + \left(\frac{n_1}{n_0(n_0-1)}+\frac{2(1-\lambda)n_0}{n_1(n_0-1)}\right) \vhn^0\right].
    \end{align*}
    }
    Here, we have introduced the hyper-parameter $\lambda \in \R$ which trades off the contributions of the two diagonal terms to the bound on the cross term. We can choose it so that the two constants in front of $\vhn^1$ and $\vhn^0$ agree. This is achieved at
    \begin{equation*}
        \lambda = \frac{1}{2} \frac{n_1^2(n_1-1) - n_0^2(n_0-1) + 2n_0^2(n_1-1)}{n_0^2(n_1-1)+n_1^2(n_0-1)}.
    \end{equation*}
    Inserting this choice of $\lambda$ into the upper bound of $L(Z)$, we arrive at the result
    \begin{equation*}
        L(Z) \leq 2 (b-a)^2 \frac{n\,(n_0^2+n_1^2)^2}{n_0n_1^3(n_0-1) + n_1 n_0^3(n_1-1)} (\vhn^1+\vhn^0).
    \end{equation*}
\end{proof}

\end{document}
